\documentclass[11pt,a4paper]{article}
\usepackage[utf8]{inputenc}
\usepackage[T1]{fontenc}
\usepackage{lmodern}
\IfFileExists{french.ldf}{%
  \usepackage[french]{babel}%
}{%
  \usepackage[english]{babel}%
  \addto\captionsenglish{%
    \renewcommand{\contentsname}{Table des matières}%
    \renewcommand{\listfigurename}{Table des figures}%
    \renewcommand{\listtablename}{Liste des tableaux}%
    \renewcommand{\chaptername}{Chapitre}%
    \renewcommand{\figurename}{Figure}%
    \renewcommand{\tablename}{Tableau}%
    \renewcommand{\bibname}{Bibliographie}%
    \renewcommand{\appendixname}{Annexe}}%
}
\usepackage[a4paper,margin=2.4cm,headheight=14pt]{geometry}
\usepackage{amsmath,amssymb,amsthm,mathtools}
\usepackage{graphicx}
\usepackage{booktabs,tabularx,array,multirow}
\usepackage[table]{xcolor}
\usepackage{tikz}
\usetikzlibrary{arrows.meta,positioning,shapes.geometric,calc,babel}
\usepackage[most]{tcolorbox}
\usepackage{listings}
\usepackage{caption}
\usepackage{float}
\usepackage{enumitem}
\usepackage{fancyhdr}
\usepackage{titlesec}
\usepackage[colorlinks=true,linkcolor=bleu,urlcolor=bleu,citecolor=bleu]{hyperref}

\definecolor{bleu}{RGB}{20,60,130}
\definecolor{gris}{RGB}{245,246,248}
\definecolor{vert}{RGB}{30,120,60}
\definecolor{orange}{RGB}{200,100,20}

\graphicspath{{../figures/}{figures/}}
\setlength{\parindent}{0pt}
\setlength{\parskip}{5pt}
\renewcommand{\arraystretch}{1.25}
\captionsetup{font=small,labelfont=bf}
\newcolumntype{L}{>{\raggedright\arraybackslash}X}

% En-têtes
\pagestyle{fancy}
\fancyhf{}
\fancyhead[L]{\small\nouppercase{\leftmark}}
\fancyfoot[R]{\thepage}
\renewcommand{\headrulewidth}{0.4pt}
\fancypagestyle{plain}{\fancyhf{}\fancyfoot[R]{\thepage}\renewcommand{\headrulewidth}{0pt}}

\ifdefined\chapter
\titleformat{\chapter}[display]{\normalfont\bfseries\color{bleu}}{\large\chaptertitlename\ \thechapter}{6pt}{\LARGE}
\titlespacing*{\chapter}{0pt}{-10pt}{18pt}
\fi
\titleformat{\section}{\normalfont\Large\bfseries\color{bleu}}{\thesection}{0.6em}{}
\titleformat{\subsection}{\normalfont\large\bfseries}{\thesubsection}{0.6em}{}

% Encadrés
\newtcolorbox{demo}[1][]{colback=gris,colframe=bleu!60,boxrule=0.5pt,arc=2pt,left=6pt,right=6pt,top=4pt,bottom=4pt,title={\textbf{Démonstration#1}},fonttitle=\small,coltitle=white,colbacktitle=bleu!70}
\newtcolorbox{retenir}{colback=vert!5,colframe=vert!70,boxrule=0.5pt,arc=2pt,left=6pt,right=6pt,top=4pt,bottom=4pt,title={\textbf{À retenir}},fonttitle=\small,coltitle=white,colbacktitle=vert!70}
\newtcolorbox{question}[1]{colback=orange!5,colframe=orange!80,boxrule=0.5pt,arc=2pt,left=6pt,right=6pt,top=4pt,bottom=4pt,title={\textbf{Question : #1}},fonttitle=\small,coltitle=white,colbacktitle=orange!80}

% Code
\lstset{language=Python,basicstyle=\ttfamily\footnotesize,keywordstyle=\color{bleu}\bfseries,
  commentstyle=\color{vert}\itshape,stringstyle=\color{orange},numbers=left,numberstyle=\tiny\color{gray},
  numbersep=6pt,frame=single,rulecolor=\color{gray!40},backgroundcolor=\color{gris},breaklines=true,
  showstringspaces=false,tabsize=4,columns=fullflexible,keepspaces=true,
  literate={é}{{\'e}}1 {è}{{\`e}}1 {à}{{\`a}}1 {ù}{{\`u}}1 {ê}{{\^e}}1 {ô}{{\^o}}1 {û}{{\^u}}1 {ç}{{\c c}}1
           {î}{{\^i}}1 {â}{{\^a}}1 {É}{{\'E}}1 {À}{{\`A}}1 {σ}{{$\sigma$}}1 {²}{{$^2$}}1 {Δ}{{$\Delta$}}1
           {−}{{-}}1 {√}{{$\surd$}}1 {μ}{{$\mu$}}1 {≈}{{$\approx$}}1 {∈}{{$\in$}}1 {±}{{$\pm$}}1 {→}{{$\to$}}1 {«}{{\guillemotleft}}1 {»}{{\guillemotright}}1}

% Raccourcis
\newcommand{\E}{\mathbb{E}}
\newcommand{\Q}{\mathbb{Q}}
\newcommand{\Var}{\operatorname{Var}}
\newcommand{\Cov}{\operatorname{Cov}}
\newcommand{\Corr}{\operatorname{Corr}}
\newcommand{\Ncal}{\mathcal{N}}
\newcommand{\1}{\mathbf{1}}
\newcommand{\CBS}{C_{BS}}
\newcommand{\vir}{{,}}

\geometry{margin=2cm}
\graphicspath{{slides_img/}{../figures/}}
\pagestyle{fancy}\fancyhf{}
\fancyhead[L]{\small Chra7 el présentation b tounsi — Monte-Carlo, call européen}\fancyhead[R]{\small\thepage}
\titleformat{\section}{\normalfont\large\bfseries\color{bleu}}{}{0em}{}
\setcounter{secnumdepth}{0}

\newcommand{\chnowa}[1]{\par\noindent\textbf{\color{bleu}Chnowa 3malna ?}\enspace #1\par}
\newcommand{\aalech}[1]{\par\noindent\textbf{\color{orange}3lech ?}\enspace #1\par}
\newcommand{\kifech}[1]{\par\noindent\textbf{\color{bleu}Kifech ?}\enspace #1\par}
\newcommand{\natija}[1]{\par\noindent\textbf{\color{vert}El natija :}\enspace #1\par}
\newcommand{\aqoul}[1]{\begin{tcolorbox}[colback=orange!6,colframe=orange!80,boxrule=0.4pt,arc=2pt,left=5pt,right=5pt,top=2pt,bottom=2pt]\small\textbf{9oulha haka lel prof :} « #1 »\end{tcolorbox}}
\newcommand{\slide}[3]{\clearpage\section{Slide #1 — #2}\begin{center}\setlength{\fboxsep}{0pt}\fbox{\includegraphics[width=0.66\textwidth]{#3}}\end{center}\vspace{-4pt}}

\begin{document}

\begin{center}
{\LARGE\bfseries\color{bleu} Chra7 el présentation, slide b slide}\\[4pt]
{\large b tounsi, b kelma sahla}\\[2pt]
{\small Projet Math et Assurance — ESPRIT 2026–2027 — M. Chebbi}
\end{center}

\begin{tcolorbox}[colback=gris,colframe=bleu!50,boxrule=0.4pt]
\small Kol slide fiha : \textbf{chnowa 3malna}, \textbf{3lech}, \textbf{kifech}, w \textbf{el natija}. W fi le5er, jomla sghira bel français \textbf{t9oulha lel prof} kif ma hiya.
\end{tcolorbox}

\subsection*{El fekra el kbira (fi 4 jomal)}
\begin{enumerate}[nosep]
\item El call européen 3andou \textbf{prix exact} : $C = 10{,}4506$ (Black-Scholes). Hadha el « jaweb el s7i7 ».
\item 7sebna nafs el prix b \textbf{Monte-Carlo} : n3mlou moyenne mta3 1 million payoff simulé. Les tirages $Z$ ysan3hom \textbf{Box-Muller}.
\item Chefna elli Monte-Carlo yjib el jaweb el s7i7 (dakhel l'IC à 95 \%).
\item \textbf{7assanna} b zouz méthodes (antithétique w contrôle) bech ywalli as7 b nafs 3add el simulations.
\end{enumerate}
\textbf{Les paramètres :} $S_0 = 100$, $K = 100$, $T = 1$ an, $r = 5\,\%$, $\sigma = 20\,\%$, $N = 10^6$.

%=====================================================================
\slide{1}{El titre}{s01.jpg}
\chnowa{Bech nwarriw el code mta3na partie b partie : el code w chnowa yokhrej.}
\aalech{El chiffres lkol elli fel rapport jeyin mel code. Ki nwarriw el code, el prof ynajjem yet2akked men kol ra9m.}
\aqoul{Nous allons présenter le notebook partie par partie : le code et ce qu'il affiche.}

%=====================================================================
\slide{2}{El plan : 9 parties}{s02.jpg}
\chnowa{El code 9assamnah 3la 9 parties sghar, kol wa7da fiha fekra wa7da.}
\aalech{Hadha el tartib el manti9i : n7adhrou les paramètres w Box-Muller, n7esbou el prix el s7i7, ba3d Monte-Carlo, ba3d n7assnouh, w fel le5er n9arnou.}
\aqoul{Le code est court : des variables simples, et les formules écrites directement dans les fonctions.}

%=====================================================================
\slide{3}{Partie 1 : imports w paramètres}{s03.jpg}
\chnowa{Ktebna les paramètres marra wa7da, b variables 3adiyin ($S_0$, $K$, $T$, $r$, $\sigma$, $N$), w fixina el générateur b graine 2026.}
\aalech{El ordinateur ma ya3mlch hasard 7a9i9i : ya7seb ar9am men graine. Nafs el graine = nafs el ar9am. Donc ki el prof y3awed y5addem el code, yal9a \textbf{nafs el résultats bel dhabt}.}
\kifech{\texttt{rng = np.random.default\_rng(2026)}. W \texttt{z95 = 1.96} houwa el ra9m elli 95 \% mel loi normale bin $-1{,}96$ w $1{,}96$.}
\aqoul{Le générateur est fixé : nos résultats sont reproductibles à l'identique.}

%=====================================================================
\slide{4}{Partie 2 : Box-Muller}{s04.jpg}
\chnowa{7awwelna les uniformes mta3 el générateur l gaussiennes $Z\sim\Ncal(0,1)$ b \textbf{la méthode de Box-Muller}, b \textbf{Z wa7da} kol marra.}
\aalech{El générateur ma yaatich ken ar9am uniformes bin 0 w 1. Ama $S_T$ yest7a99 gaussiennes. Box-Muller formule \textbf{exacte} (ma fihech approximation) w sahla.}
\kifech{$Z = \sqrt{-2\ln U_1}\,\cos(2\pi U_2)$.\\
\emph{El démo b ikhtisar :} $R = \sqrt{-2\ln U_1}$ 3andou densité $r\,e^{-r^2/2}$, w $\Theta = 2\pi U_2$ uniforme. Ki net3addaw lel coordonnées polaires, la densité twalli $\varphi(x)\varphi(y)$, donc $R\cos\Theta$ tetba3 $\Ncal(0,1)$.\\
\emph{Détail :} nekhdhou $1-U$ bech ma n7esbouch $\ln 0$.}
\natija{3la 1 million tirage : moyenne $-0{,}0010$ ($\approx0$), variance $0{,}9987$ ($\approx1$), $\Q(|Z|\le1{,}96) = 0{,}9503$ ($\approx0{,}95$). El histogramme lasse9 3al courbe.}
\aqoul{Box-Muller transforme deux uniformes en une gaussienne, de façon exacte ; nos vérifications confirment la loi N(0,1).}

%=====================================================================
\slide{5}{Partie 3 : el prix exact (Black-Scholes)}{s05.jpg}
\chnowa{Programmina la formule $C = S_0N(d_1) - Ke^{-rT}N(d_2)$.}
\aalech{Hiya \textbf{el référence} : bla biha ma na3rfouch 9addech Monte-Carlo ghalet.}
\kifech{\texttt{norm.cdf} hiya $N(x)$. $d_1 = 0{,}35$, $d_2 = 0{,}15$, $N(d_1) = 0{,}6368$, $N(d_2) = 0{,}5596$.}
\natija{$C = 10{,}4506$.}
\aqoul{Le prix exact 10,4506 est notre référence pour mesurer l'erreur de Monte-Carlo.}

%=====================================================================
\slide{6}{Partie 4 : Monte-Carlo standard}{s06.jpg}
\chnowa{7sebna el prix b moyenne mta3 1 million payoff simulé, m3a l'erreur standard $SE$ w intervalle de confiance. El formule mta3 $S_T$ w el payoff ktebnahom \textbf{toul dakhel el fonction} (bla def).}
\aalech{El prix houwa espérance. \textbf{Loi des grands nombres} : el moyenne tmchi lel prix el s7i7. \textbf{TCL} : yaatina el marge d'erreur.}
\kifech{4 étapes : (1) $Z$ b Box-Muller ; (2) $S_T = S_0e^{(r-\sigma^2/2)T+\sigma\sqrt TZ}$ (solution exacte, ma fammech erreur de discrétisation) ; (3) payoff actualisé $Y = e^{-rT}\max(S_T-K,0)$ ; (4) moyenne. IC $=$ prix $\pm1{,}96\,SE$.}
\natija{$\hat C = 10{,}4370$, IC $= [10{,}4081\ ;\ 10{,}4658]$ : el prix exact dakhel l'IC. \textbf{Erreur} $|C_{\text{analytique}} - \hat C| = 0{,}013624$ (nafs el 7aja elli 3malneha m3a $c$ : el analytique naa9es mta3na, en valeur absolue).}
\aqoul{La moyenne des payoffs simulés converge vers le prix ; l'IC à 95 \% contient le prix exact.}

%=====================================================================
\slide{7}{L'IC 7asb $N$}{s07.jpg}
\chnowa{3awedna el calcul l 25 valeur mta3 $N$, men 1\,000 l 1\,000\,000, w rsamna el prix w l'IC (kif ma yotlob l'énoncé).}
\natija{L'IC ydhaya9 3la 10,4506, w el prix exact dakhel l'IC fel 25 kol.}
\aqoul{L'intervalle se resserre autour du prix exact quand N augmente.}

%=====================================================================
\slide{8}{Les puissances de $N$}{s08.jpg}
\chnowa{7sebna el prix l $N = 10, 100, 1\,000, \dots, 10^6$, m3a l'IC w el ghalta $|\hat C - C|$, fi \textbf{tableau} (bla graphique).}
\aalech{Bech nchoufou kifech el précision tzid ki $N$ yetdhreb fi 10.}
\natija{Kol ma $N$ yetdhreb fi 10, l'IC yet9assem 3la $\sqrt{10}\approx3{,}16$. B $N = 10^6$ : $10{,}4371$, ghalta $0{,}0134$. El prix exact dakhel l'IC dima.}
\kifech{\textbf{Ken yes2lek :} « 3lech el ghalta b $10^4$ akbar men $10^3$ ? » $\to$ el ghalta mta3 simulation wa7da hasard ; elli dima yonzel houwa \textbf{l'IC}.}
\aqoul{À chaque puissance de 10, l'intervalle est divisé par environ 3,16 : l'erreur décroît en $1/\sqrt N$.}

%=====================================================================
\slide{9}{La trajectoire : stabilité}{s09.jpg}
\chnowa{Tba3na el moyenne $\hat C_n$ mta3 el $n$ payoffs el lowlin, men $n = 1$ l 1 million, 3la simulation wa7da.}
\aalech{Bech nwarriw elli Monte-Carlo \textbf{yestabilisa} : kol ma nzidou simulations, el moyenne t9ell ma tet7arrek.}
\kifech{\texttt{np.cumsum(Y) / n}. El pointillés : l'enveloppe $C\pm1{,}96\,s/\sqrt n$.}
\natija{Fel bdeya tet7arrek barcha, ba3d testabilisa : $10{,}6402$ b 1\,000, $10{,}4308$ b 1 million, dakhel l'enveloppe.}
\aqoul{La trajectoire se stabilise autour du prix exact : c'est la loi des grands nombres.}

%=====================================================================
\slide{10}{Partie 5 : l'antithétique}{s10.jpg}
\chnowa{Kol $Z$, nesta3mlou zeda $-Z$, w na3mlou moyenne mta3 el zouz payoffs. B nafs el budget : $N/2$ tirage $=$ $N$ payoff.}
\aalech{Bech n9allou el variance bla ma nzidou simulations. $Z$ w $-Z$ 3andhom nafs el loi, donc ma fammech biais.}
\kifech{Gain $= 1/(1+\rho)$, avec $\rho = \Corr(Y^+,Y^-)$. Fel code, $-Z$ ywalli $-\sigma\sqrt T Z$ fel formule mta3 $S_T$.}
\natija{$\hat C_a = 10{,}4600$, $SE = 0{,}01039$ (contre $0{,}01471$) : gain $= 2{,}00$, donc $\rho\approx-0{,}5$. \textbf{Erreur} $|C - \hat C_a| = 0{,}009445$. Mouch akthar khater el payoff \textbf{convexe}.}
\aqoul{L'antithétique divise la variance par 2 ; le gain est limité car le payoff du call est convexe.}

%=====================================================================
\slide{11}{Partie 6 : variable de contrôle w erreur 3la $c$}{s11.jpg}
\chnowa{Sal7na el payoff b $X = S_T$ : $Y_c = Y - c^*(X - \E X)$. W 9arna $c^*$ mta3na m3a \textbf{$c$ analytique} (el formule exacte).}
\aalech{$Y$ w $S_T$ mertbtin barcha : n9ass9ou el ziyada elli na3rfouha. W el erreur 3la $c$ twarri elli el calcul mta3na s7i7.}
\kifech{\textbf{$c^*$ :} el variance $g(c)$ parabole, ndériviw : $c^* = \Cov(Y,X)/\Var(X)$.\\
\textbf{$c$ analytique :} b la loi log-normale : $\E[S_T\1_{S_T>K}] = S_0e^{rT}N(d_1)$, $\E[S_T^2\1_{S_T>K}] = S_0^2e^{(2r+\sigma^2)T}N(d_1+\sigma\sqrt T)$, $\Var S_T = S_0^2e^{2rT}(e^{\sigma^2T}-1)$, w ba3d $c = \Cov/\Var$.}
\natija{$c$ analytique $= 0{,}640762$, $c^* = 0{,}640505$ : \textbf{erreur $= 0{,}000257$} (0,04 \%). $\hat C_c = 10{,}4472$, $SE = 0{,}00562$, gain $= 6{,}86$, \textbf{erreur sur le prix} $|C - \hat C_c| = 0{,}003408$.}
\aqoul{Notre c simulé (0,6405) est très proche du c analytique (0,6408) : l'erreur n'est que de 0,00026.}

%=====================================================================
\slide{12}{Partie 8 : el comparaison}{s12.jpg}
\chnowa{7attina el 3 méthodes fi tableau wa7ed, b nafs el budget (1 million).}
\kifech{MSE $= SE^2$ (ma fammech biais). Gain $= (SE_{\text{standard}}/SE)^2$. Efficacité $=$ gain $\times$ (temps standard / temps). El temps n7esbouh \textbf{moyenne 3la 5 exécutions} (mesure wa7da fiha barcha hasard), b générateur wa7dou bech ma yetbaddlouch el résultats.}
\natija{Antithétique : gain $\times2{,}00$, efficacité $\times3{,}5$. Contrôle : gain $\times6{,}86$, efficacité $\times4{,}7$ : \textbf{el a7sen}. El 3 IC fihom el prix exact.}
\aqoul{Le meilleur estimateur est la variable de contrôle : on le choisit sur l'erreur standard, pas sur l'écart d'une seule simulation.}

%=====================================================================
\slide{13}{La courbe de comparaison des 3 méthodes}{s13.jpg}
\chnowa{Rsamna el trajectoire mta3 el 3 méthodes 3la nafs el graphique.}
\aalech{El tableau yaatina el précision b 1 million bark ; el courbe twarri kifech kol méthode t9arreb lel prix l kol $N$.}
\natija{El 3 courbes yousslou l 10,4506. El contrôle (el a5dher) \textbf{y9ell ma yet7arrek} : houwa el as7.}
\aqoul{Les trois méthodes convergent vers le prix exact ; la variable de contrôle oscille le moins.}

%=====================================================================
\slide{14}{Les puissances de 10 : el 3 méthodes}{s14.jpg}
\chnowa{3awedna el 7keya mta3 slide 8 ($N = 10, 100, \dots, 10^6$) ama bel 3 méthodes, w 7sebna l kol wa7da el erreur $|C_{\text{analytique}} - \hat C|$.}
\aalech{Bech nchoufou el 3 méthodes kifech yt9arbou lel prix ki $N$ ykbar.}
\natija{El prix yet9arreb l 10,4506 fel 3 (tableau bark, bla graphique). B $N = 10^6$ : erreur $0{,}0020$ (standard), $0{,}0187$ (antithétique), $0{,}0010$ (contrôle).}
\kifech{\textbf{Ken yes2lek :} « 3lech b $N = 10$ l'antithétique 3tat 4,70 ? » $\to$ 5 paires bark, el TCL mazel ma yekhdemch. « 3lech b $10^3$ el contrôle ghalet akthar ? » $\to$ el erreur mta3 simulation wa7da hasard ; bech na3rfou chkoun el as7, n3awdou el calcul barcha marrat b $N$ fixe : slide 15.}
\aqoul{Pour les trois méthodes, le prix se rapproche du prix exact quand N augmente.}

%=====================================================================
\slide{15}{Stabilité w erreur moyenne b $N = 10^6$ fixe}{s15.jpg}
\chnowa{Thabbatna $N = 10^6$ w 3awedna el calcul \textbf{50 marra} b kol méthode (tirages jdod kol marra). Rsamna \textbf{courbe wa7da} : kol méthode 3andha colonne fiha el 50 prix mte3ha, w el ligne el 7amra hiya el prix exact.}
\aalech{Simulation wa7da tnajjem tjina 9rib mel prix b el hasard. Ki n3awdou 50 marra, nchoufou chkoun el méthode elli \textbf{testabilisa} akther : el colonne elli \textbf{adhya9} hiya el as7.}
\kifech{\textbf{Erreur moyenne} $=$ moyenne mta3 el 50 erreurs $|C_{\text{analytique}} - \hat C|$. Fel code : \texttt{np.abs(C\_exact - prix\_M).mean()}. Maktouba fel légende mta3 el courbe w fel tableau.}
\natija{Erreur moyenne : \textbf{0,0127} (standard), \textbf{0,0080} (antithétique), \textbf{0,0044} (contrôle). El contrôle a9al b 3 marrat mel standard. El moyennes el 3 9rab mel 10,4506 : ma fammech biais.}
\kifech{\textbf{Ken yes2lek :} « chnowa el far9 bin erreur moyenne w écart-type ? » $\to$ el erreur moyenne : 9addech nab3dou 3al \textbf{prix exact} ; el écart-type : 9addech el 50 prix mfarr9in 3la \textbf{moyennethom}. Khater ma fammech biais, el zouz y3tiw nafs el tartib.}
\aqoul{À N fixé, la variable de contrôle est la plus stable : son erreur moyenne est environ 3 fois plus petite que celle du standard.}

%=====================================================================
\slide{16}{Partie 9 : la courbe de gain}{s16.jpg}
\chnowa{Rsamna el $SE$ mta3 el 3 méthodes 7asb $N$, w el gain 7asb $N$.}
\aalech{Bech nwarriw \textbf{kifech} t5addem la réduction de variance, w 3lech l'antithétique hiya l a9al.}
\natija{\textbf{Ysar :} 3 droites parallèles (pente $-1/2$) : el sor3a ma tetbaddelch, ama el droite tahbet. \textbf{Ymin :} el gain sabet ($\approx2$ w $\approx7$). \textbf{3lech l'antithétique l a9al :} $1/(1+\rho)$ m3a $\rho\approx-0{,}5$ ya3ti 2 ; el contrôle $1/(1-\rho^2)$ m3a $\rho\approx0{,}92$ ya3ti 7.}
\aqoul{La réduction de variance ne change pas la vitesse en $1/\sqrt N$, elle abaisse la droite ; et le gain ne dépend pas de N.}

%=====================================================================
\slide{17}{La conclusion}{s17.jpg}
\natija{(1) Box-Muller ysan3 gaussiennes s7a7. (2) Monte-Carlo yjib el prix el s7i7 10,4506 dakhel l'IC, w yestabilisa. (3) La variable de contrôle t9assem el variance 3la 6,9, w el $c$ mta3na ghalet b 0,00026 bark.}
\aqoul{Avec Box-Muller, Monte-Carlo retrouve le prix exact, et la variable de contrôle divise la variance par 6,9.}

\end{document}
